\exercisesection{\S~13}

\begin{enumerate}
\def\labelenumi{\arabic{enumi})}
\tightlist
\item
  The dimensions \(\leq 80\) of splittable simple Lie algebras are:
\end{enumerate}

3 (A \(_{1}\) = B \(_{1}\) = C \(_{1}\) ), 8 (A \(_{2}\) ), 10 (B \(_{2}\) = C \(_{2}\) ), 14 (G \(_{2}\) ), 15 (A \(_{3}\) = D \(_{3}\) ), 21 (B \(_{3}\) and C \(_{3}\) ), 24 (A \(_{4}\) ), 28 (D \(_{4}\) ), 35 (A \(_{5}\) ), 36 (B \(_{4}\) and C \(_{4}\) ), 45 (D \(_{5}\) ), 48 (A \(_{6}\) ), 52 (F \(_{4}\) ), 55 (B \(_{5}\) and C \(_{5}\) ), 63 (A \(_{7}\) ), 66 (D \(_{6}\) ), 78 (B \(_{6}\) , C \(_{6}\) and E \(_{6}\) ), 80 (A \(_{8}\) ).

\begin{enumerate}
\def\labelenumi{\arabic{enumi})}
\setcounter{enumi}{1}
\tightlist
\item
  Let \((\mathfrak{g},\mathfrak{h})\) be a split simple Lie algebra, and B a basis of \(R(\mathfrak{g},\mathfrak{h})\) . The simple \(\mathfrak{g}\) -modules \(E(\lambda)\) , \(\lambda \in P_{++} - \{0\}\) , of minimum dimension are those for which \(\lambda\) is one of the following weights:
\end{enumerate}

\(\varpi_{1}\) (A \(_{1}\) ); \(\varpi_{1}\) and \(\varpi_{l}\) (A \(_{l}\) , \(l \geq 2\) ); \(\varpi_{1}\) (B \(_{l}\) and C \(_{l}\) , \(l \geq 2\) ); \(\varpi_{1}\) , \(\varpi_{3}\) and \(\varpi_{4}\) (D \(_{4}\) ); \(\varpi_{1}\) (D \(_{l}\) , \(l \geq 5\) ); \(\varpi_{1}\) and \(\varpi_{6}\) (E \(_{6}\) ); \(\varpi_{7}\) (E \(_{7}\) ); \(\varpi_{8}\) (E \(_{8}\) ); \(\varpi_{4}\) (F \(_{4}\) ); \(\varpi_{1}\) (G \(_{2}\) ). Any two such modules can be transformed into each other by an automorphism of g.

Type \(E_{8}\) is the only one for which the adjoint representation is of minimum dimension.

\begin{enumerate}
\def\labelenumi{\arabic{enumi})}
\setcounter{enumi}{2}
\tightlist
\item
  \begin{enumerate}
  \def\labelenumii{\alph{enumii})}
  \tightlist
  \item
    Define an isomorphism from \(\mathfrak{sl}(4,k)\) to the orthogonal algebra \(\mathfrak{o}_S(6,k)\) . (Use the fact that the representation \(\bigwedge^2\sigma\) of no. 1.V is orthogonal and of dimension 6.) The two types of irreducible representation of \(\mathfrak{sl}(4,k)\) of degree 4 correspond to the two semi-spinor representations of \(\mathfrak{o}_S(6,k)\) .
  \end{enumerate}
\end{enumerate}

\begin{enumerate}
\def\labelenumi{\alph{enumi})}
\setcounter{enumi}{1}
\item
  Define an isomorphism from \(\mathfrak{sp}(4,k)\) to the orthogonal algebra \(\mathfrak{o}_S(5,k)\) . (Use the fact that the representation \(\sigma_2\) of no. 3.V is orthogonal and of dimension 5.) The irreducible representation of \(\mathfrak{sp}(4,k)\) of degree 4 corresponds to the spinor representation of \(\mathfrak{o}_S(5,k)\) .
\item
  Define an isomorphism \(\mathfrak{sl}(2,k)\times\mathfrak{sl}(2,k)\to\mathfrak{o}_{S}(4,k)\) by using the tensor product of the identity representations of the two \(\mathfrak{sl}(2,k)\) factors. Recover this result by means of Chap. I, §6, Exerc. 26.
\end{enumerate}

\begin{enumerate}
\def\labelenumi{\arabic{enumi})}
\setcounter{enumi}{3}
\tightlist
\item
  Let \(S\) be the square matrix of order \(n\)
\end{enumerate}

\[
(\delta_ {i, n + 1 - i}) = \left( \begin{array}{c c c c c} 0 & 0 & \dots & 0 & 1 \\ 0 & 0 & \dots & 1 & 0 \\ \vdots & \vdots & \ddots & \vdots & \vdots \\ 0 & 1 & \dots & 0 & 0 \\ 1 & 0 & \dots & 0 & 0 \end{array} \right).
\]

The elements of \(\mathfrak{o}_{S}(n,k)\) are the matrices \((a_{ij})\) that are anti-symmetric with respect to the second diagonal:

\[
a _ {i, j} = - a _ {n + 1 - j, n + 1 - i} \quad \text {   for   every   pair   } (i, j).
\]

The algebra \(\mathfrak{o}_S(n,k)\) is splittable simple of type \(\mathrm{D}_{n/2}\) if \(n\) is even and \(\geq 6\) , and of type \(\mathrm{B}_{(n-1)/2}\) if \(n\) is odd and \(\geq 5\) . The diagonal (resp. upper triangular) elements of \(\mathfrak{o}_S(n,k)\) form a splitting Cartan (resp. Borel) subalgebra.

\begin{enumerate}
\def\labelenumi{\arabic{enumi})}
\setcounter{enumi}{4}
\item
  The notations are those of no. 1.IV, type \(\mathrm{A}_l\) . Show that, if \(n\) is \(\geq 0\) , \(\mathbf{S}^n\sigma\) is an irreducible representation of \(\mathfrak{sl}(l + 1,k)\) of highest weight \(n\varpi_1\) . (Realise \(\mathbf{S}^n\sigma\) on the space of homogeneous polynomials of degree \(n\) in \(\mathrm{X}_0,\dots ,\mathrm{X}_l\) , and observe that, up to homothety, the only polynomial \(f\) such that \(\partial f / \partial \mathrm{X}_i = 0\) for \(i\geq 1\) is \(\mathrm{X}_0^n\) .) Show that all the weights of this representation are of multiplicity 1.
\item
  The notations are those of no. 2.IV, type \(\mathrm{B}_l\) . Show that, if \(1 \leq r \leq l - 1\) , the dimension of \(\operatorname{E}(\varpi_r)\) is \(\binom{2l + 1}{r}\) , and deduce another proof of the fact that \(\bigwedge^r \sigma\) is a fundamental representation of highest weight \(\varpi_r\) .
\item
  The notations are those of no. 2.VII, type \(\mathrm{B}_l\) . Show that \(\mathbf{O}_0^+ (\Psi)\) is the commutator subgroup of \(\mathbf{SO}(\Psi)\) . (Remark that \(\mathbf{O}(\Psi)\) is equal to \(\{\pm 1\} \times \mathbf{SO}(\Psi)\) , hence has the same commutator subgroup as \(\mathbf{SO}(\Psi)\) , and apply Algebra, Chap. IX, §9, Exerc. 11 b).) Deduce that \(\operatorname{Aut}_e(\mathfrak{g}) = \mathbf{O}_0^+ (\Psi)\) .
\item
  The notations are those of no. 3.IV, type \(C_{l}\) ( \(l \geq 1\) ). Show that \(S^{2}\sigma\) is equivalent to the adjoint representation of g.
\item
  The notations are those of no. 3.VII, type \(C_l\) ( \(l \geq 1\) ). In particular, we identify \(\mathrm{Aut}_e(\mathfrak{g})\) with a subgroup of \(\mathbf{Sp}(\varPsi)/\{\pm 1\}\) . Show that the image in \(\mathbf{Sp}(\varPsi)/\{\pm 1\}\) of a symplectic transvection (Algebra, Chap. IX, §4, Exerc. 6) belongs to \(\mathrm{Aut}_e(\mathfrak{g})\) . Deduce that \(\mathrm{Aut}_e(\mathfrak{g}) = \mathbf{Sp}(\varPsi)/\{\pm 1\}\) (Algebra, Chap. IX, §5, Exerc. 11), and that \(\mathrm{Aut}(\mathfrak{g})/\mathrm{Aut}_e(\mathfrak{g})\) can be identified with \(k^*/k^{*2}\) .
\end{enumerate}

¶ 10) The notations are those of no. 4.IV, type D \(_{l}\) (l ≥ 2).

\begin{enumerate}
\def\labelenumi{\alph{enumi})}
\tightlist
\item
  Let \(x\) and \(y\) be the elements of \(\bigwedge^{l} \mathrm{V}\) defined by
\end{enumerate}

\[
x = e _ {1} \wedge \dots \wedge e _ {l - 1} \wedge e _ {l} \quad \text { and } \quad y = e _ {1} \wedge \dots \wedge e _ {l - 1} \wedge e _ {- l}.
\]

The element x is primitive of weight \(2\varpi_{l}\) and y is primitive of weight \(2\varpi_{l-1}\) . The submodule X (resp. Y) of \(\bigwedge^{l}V\) generated by x (resp. y) is isomorphic to \(\mathrm{E}(2\varpi_{l})\) (resp. \(\mathrm{E}(2\varpi_{l-1})\) ). Show, by calculating dimensions, that \(\bigwedge^{l}V = X \oplus Y\) ; in particular, \(\bigwedge^{l}V\) is the sum of two non-isomorphic simple modules.

\begin{enumerate}
\def\labelenumi{\alph{enumi})}
\setcounter{enumi}{1}
\tightlist
\item
  Let \(e = e_{1} \wedge \cdots \wedge e_{l} \wedge e_{-1} \wedge \cdots \wedge e_{-l} \in \bigwedge^{2l}V\) , and \(\varPsi_{l}\) the extension of \(\varPsi\) to \(\bigwedge^{l}V\) . Let \(z \in \bigwedge^{l}V\) . Prove the equivalences:
\end{enumerate}

\[
z \in X \Longleftrightarrow z \wedge t = \Psi_ {l} (z, t) e \quad \text {   for   all   } t \in \bigwedge^ {l} V
\]

\[
z \in \mathrm{Y} \Longleftrightarrow z \wedge t = - \Psi_ {l} (z, t) e \quad \text {   for   all   } t \in \bigwedge^ {l} \mathrm{V}.
\]

(If \(X'\) and \(Y'\) denote the subspaces defined by the right-hand sides, prove first that \(X'\) and \(Y'\) are stable under \(\mathfrak{g}\) and contain \(x\) and \(y\) , respectively.)

\begin{enumerate}
\def\labelenumi{\alph{enumi})}
\setcounter{enumi}{2}
\item
  Assume that z is pure (Algebra, Chap. III, §11, no. 13), and denote by \(M_{z}\) the l-dimensional subspace of V associated to it. Show that \(M_{z}\) is totally isotropic if and only if z belongs to X or to Y. (Use the fact that, when \(M_{z}\) is totally isotropic, there exists an orthogonal transformation that transforms it into \(M_{x}\) ; when \(M_{z}\) is not totally isotropic, construct an l-vector t such that \(z \wedge t = 0, \Psi_{l}(z, t) = 1\) , and apply b) above.) When \(z \in X\) (resp. \(z \in Y\) ), the dimension of \(\mathrm{M}_{x}/(\mathrm{M}_{x} \cap \mathrm{M}_{z})\) is an even (resp. odd) integer, cf.~Algebra, Chap. IX, §6, Exerc. 18 d).
\item
  Let s be a direct (resp. inverse) similarity of V. Show that \(\bigwedge^{l}s\) leaves X and Y stable (resp. interchanges X and Y).
\end{enumerate}

\begin{enumerate}
\def\labelenumi{\arabic{enumi})}
\setcounter{enumi}{10}
\item
  The notations are those of no. 4.VII, type \(\mathrm{D}_l\) ( \(l \geq 3\) ). Show that \(\mathbf{O}_0^+(\varPsi)\) is the commutator subgroup of \(\mathbf{SO}(\varPsi)\) . (Apply Algebra, Chap. IX, §6, Exerc. 17 b) and Algebra, Chap. IX, §9, Exerc. 11 b). Deduce that \(\mathrm{Aut}_e(\mathfrak{g})\) is equal to the image of \(\mathbf{O}_0^+(\varPsi)\) in \(\mathbf{SO}(\varPsi)/\{\pm 1\}\) . Moreover, \(-1\) belongs to \(\mathbf{O}_0^+(\varPsi)\) if and only if \(l\) is even or \(-1\) is a square in \(k\) (Algebra, Chap. IX, §9, Exerc. 11 c)).
\item
  The Killing form of \(\mathfrak{sl}(n,k)\) is \((X,Y)\mapsto 2n\operatorname {Tr}(XY)\) . That of \(\mathfrak{sp}(n,k)\) , \(n\) even, is \((X,Y)\mapsto (n + 2)\operatorname {Tr}(XY)\) . That of \(\mathfrak{o}_S(n,k)\) , where \(S\) is non-degenerate symmetric of rank \(n\) , is \((X,Y)\mapsto (n - 2)\operatorname {Tr}(XY)\) .
\item
  The algebra of invariant polynomial functions of \(\mathfrak{g}\) is generated:
\end{enumerate}

\(a\) ) in case \(\mathrm{A}_l\) , by the functions \(X\mapsto \mathrm{Tr}(X^i)\) , \(2\leq i\leq l + 1\) ;

\begin{enumerate}
\def\labelenumi{\alph{enumi})}
\setcounter{enumi}{1}
\item
  in case \(B_{l}\) , by the functions \(X \mapsto \operatorname{Tr}(X^{2i})\) , \(1 \leq i \leq l\) ;
\item
  in case \(C_{l}\) , by the functions \(X \mapsto \operatorname{Tr}(X^{2i})\) , \(1 \leq i \leq l\) ;
\end{enumerate}

\(d\) ) in case \(\mathrm{D}_l\) , by the functions \(X \mapsto \operatorname{Tr}(X^{2i})\) , \(1 \leq i \leq l - 1\) , and by one of the two polynomial functions \(\tilde{f}\) such that \(\tilde{f}(X)^2 = (-1)^l \det(X)\) .

\begin{enumerate}
\def\labelenumi{\arabic{enumi})}
\setcounter{enumi}{13}
\tightlist
\item
  \begin{enumerate}
  \def\labelenumii{\alph{enumii})}
  \tightlist
  \item
    Let G be the group associated to \(\mathfrak{g} = \mathfrak{sl}(n, k)\) by the procedure of §7, Exerc. 26. The natural g-module structure on \(k^{n}\) gives rise to a homomorphism \(\varphi : G \to \mathbf{GL}(n, k)\) . Use loc. cit. h) to prove that \(\varphi\) is injective, and loc. cit. f) to prove that
  \end{enumerate}
\end{enumerate}

\[
\operatorname{Im} (\varphi) = \mathbf {S L} (n, k).
\]

\begin{enumerate}
\def\labelenumi{\alph{enumi})}
\setcounter{enumi}{1}
\item
  Let E be a finite dimensional \(\mathfrak{sl}(n,k)\) -module, and \(\rho\) the corresponding representation of \(\mathfrak{sl}(n,k)\) . Show that there exists a unique representation \(\pi:\mathbf{SL}(n,k)\to\mathbf{GL}(\mathrm{E})\) such that \(\pi(e^{x})=e^{\rho(x)}\) for every nilpotent element x of \(\mathfrak{sl}(n,k)\) (use a)). We say that \(\rho\) and \(\pi\) are compatible. Generalize the results proved for n=2 in §1, no. 4.
\item
  Assume that \(k\) is \(\mathbf{R}, \mathbf{C}\) or a complete field for a discrete valuation with residue field of characteristic \(\neq 0\) . Show that \(\rho\) and \(\pi\) are compatible if and only if \(\pi\) is a homomorphism of Lie groups such that \(L(\pi) = \rho\) (same method as in Exerc. 18 b) of §1).
\end{enumerate}

\(d\) ) Prove analogous results for \(\mathfrak{sp}(2n,k)\) and \(\mathbf{Sp}(2n,k)\) .

¶15) Let V be a vector space of finite dimension \(\geq 2\) , g a Lie subalgebra of End(V) and \(\theta\) an element of g. We make the following assumptions:

\begin{enumerate}
\def\labelenumi{(\roman{enumi})}
\item
  V is a semi-simple g-module;
\item
  \(\theta\) is of rank 1 (i.e.~\(\dim \operatorname{Im}(\theta) = 1\) );
\item
  the line \(\operatorname{Im}(\theta)\) generates the U(g)-module V.
\end{enumerate}

\begin{enumerate}
\def\labelenumi{\alph{enumi})}
\item
  Show that these assumptions are satisfied (for a suitable choice of \(\theta\) ) when \(\mathfrak{g} = \mathfrak{sl}(V)\) , when \(\mathfrak{g} = \mathfrak{gl}(V)\) , or when there exists a non-degenerate alternating bilinear form \(\varPsi\) on V such that \(\mathfrak{g} = \mathfrak{sp}(\varPsi)\) or \(\mathfrak{g} = k.1 \oplus \mathfrak{sp}(\varPsi)\) . In each of these cases, \(\theta\) can be taken to be a nilpotent element; in the second case (and only in that case) \(\theta\) can be taken to be a semi-simple element.
\item
  We shall prove that the four cases above are the only ones possible. We are reduced immediately to the case in which k is algebraically closed. Show that V is then a simple g-module, and that \(g = c \oplus s\) , where s is semi-simple, and c = 0 or c = k.1; show that V is not isomorphic to a tensor product of g-modules of dimension \(\geq 2\) ; deduce that s is simple.
\item
  Choose a Cartan subalgebra \(\mathfrak{h}\) of \(\mathfrak{s}\) , and a basis B of \(\mathrm{R}(\mathfrak{s},\mathfrak{h})\) . Let \(\lambda\) be the highest weight (with respect to B) of the \(\mathfrak{s}\) -module V, and let \(e\) be a non-zero element of V of weight \(\lambda\) . The highest weight of the dual module \(\mathrm{V}^*\) is \(\lambda^{*} = -w_{0}\lambda\) (§7, no. 5); let \(e^*\) be a non-zero element of \(\mathrm{V}^*\) of weight \(\lambda^*\) . Identify \(\mathrm{V} \otimes \mathrm{V}^*\) with \(\operatorname{End}(\mathrm{V})\) as usual. Show that there exist \(x \in \mathrm{V}, y \in \mathrm{V}^*\) such that \(x \otimes y \in \mathfrak{g}\) and \(\langle x, e^* \rangle \neq 0, \langle e, y \rangle \neq 0\) (take the conjugate of \(\theta\) by \(e^n\) , where \(n\) is a suitable nilpotent element of \(\mathfrak{g}\) ). Use the fact that \(\mathfrak{g}\) is an \(\mathfrak{h}\) -submodule of \(\mathrm{V} \otimes \mathrm{V}^*\) to conclude that \(\mathfrak{g}\) contains \(e \otimes e^*\) . Deduce that \(\lambda + \lambda^{*} = \tilde{\alpha}\) , where \(\tilde{\alpha}\) is the highest root of \(\mathfrak{s}\) .
\end{enumerate}

\(d\) ) Show that \(\mathfrak{s}\) is not of type \(\mathrm{B}_l\) ( \(l\geq 2\) ), \(\mathrm{D}_l\) ( \(l\geq 4\) ), \(\mathrm{E}_6,\mathrm{E}_7,\mathrm{E}_8,\mathrm{F}_4,\mathrm{G}_2\) (by Chap. VI, Tables, \(\tilde{\alpha}\) would be a fundamental weight, and hence could not be of the form \(\lambda +\lambda^{*}\) above). Deduce that \(\mathfrak{s}\) is either of type \(\mathrm{A}_l\) or of type \(\mathrm{C}_l\) , and that in the first case \(\lambda = \varpi_1\) or \(\varpi_1 = \varpi_1^*\) , and in the second case \(\tilde{\alpha} = 2\varpi_1 = \varpi_1 + \varpi_1^*\) ; since \(\mathfrak{c} = 0\) or \(k.1\) , this indeed gives the four possibilities in \(a)^{21}\) .

¶16) Let g be an absolutely simple Lie algebra of type \(A_{l}\) ( \(l \geq 2\) ), \(\bar{k}\) an algebraic closure of k, and \(\pi : \operatorname{Gal}(\bar{k}/k) \to \operatorname{Aut}(\bar{\mathbf{R}},\bar{\mathbf{B}})\) the homomorphism defined in Exerc. 8 of §5.

\begin{enumerate}
\def\labelenumi{\alph{enumi})}
\tightlist
\item
  Assume that \(\pi\) is trivial. Show that there exist exactly two two-sided ideals \(\mathfrak{m}\) and \(\mathfrak{m}'\) of \(\mathrm{U}(\mathfrak{g})\) such that \(\mathrm{D} = \mathrm{U}(\mathfrak{g}) / \mathfrak{m}\) and \(\mathrm{D}' = \mathrm{U}(\mathfrak{g}) / \mathfrak{m}'\) are central simple algebras of dimension \((l + 1)^2\) (use Exerc. 8 of §7). The principal anti-automorphism of \(\mathrm{U}(\mathfrak{g})\) interchanges \(\mathfrak{m}\) and \(\mathfrak{m}'\) ; in particular, \(\mathrm{D}'\) is isomorphic to the opposite of D. The composite \(\mathfrak{g} \to \mathrm{U}(\mathfrak{g}) \to \mathrm{D}\) identifies \(\mathfrak{g}\) with the Lie subalgebra \(\mathfrak{sl}_{\mathrm{D}}\) of D consisting of the elements of trace zero. Moreover, \(\mathfrak{g}\) is splittable (and hence isomorphic to \(\mathfrak{sl}(l + 1,k)\) ) if and only if D is isomorphic to \(\mathbf{M}_{l + 1}(k)\) .
\end{enumerate}

Conversely, if \(\Delta\) is a central simple algebra of dimension \((l + 1)^2\) , the Lie algebra \(\mathfrak{sl}_{\Delta}\) is absolutely simple of type \(A_l\) , and the corresponding homomorphism \(\pi\) is trivial. Two such algebras \(\mathfrak{sl}_{\Delta}\) and \(\mathfrak{sl}_{\Delta'}\) are isomorphic if and only if \(\Delta\) and \(\Delta'\) are isomorphic or anti-isomorphic.

\begin{enumerate}
\def\labelenumi{\alph{enumi})}
\setcounter{enumi}{1}
\tightlist
\item
  Assume that \(\pi\) is non-trivial. Since \(\operatorname{Aut}(\bar{\mathbf{R}},\bar{\mathbf{B}})\) has two elements, the kernel of \(\pi\) is an open subgroup of \(\operatorname{Gal}(\bar{k}/k)\) of index 2, which corresponds by Galois theory to a quadratic extension \(k_{1}\) of k. Denote the non-trivial involution of \(k_{1}\) by \(x\mapsto\bar{x}\) .
\end{enumerate}

Show that there exists a unique two-sided ideal \(\mathfrak{m}\) of \(\mathrm{U}(\mathfrak{g})\) such that \(\mathrm{D} = \mathrm{U}(\mathfrak{g}) / \mathfrak{m}\) is a simple algebra of dimension \(2(l + 1)^2\) whose centre is a quadratic extension of \(k\) (same method). The centre of D can be identified with \(k_{1}\) . The ideal \(\mathfrak{m}\) is stable under the principal anti-automorphism of \(\mathrm{U}(\mathfrak{g})\) ; that defines by passage to the quotient an involutive anti-automorphism \(\sigma\) of D such that \(\sigma(x) = \bar{x}\) for all \(x \in k_{1}\) . The composite map \(\mathfrak{g} \to \mathrm{U}(\mathfrak{g}) \to \mathrm{D}\) identifies \(\mathfrak{g}\) with the Lie subalgebra \(\mathfrak{s}\mathfrak{u}_{\mathrm{D},\sigma}\) of D consisting of the elements \(x\) such that \(\sigma(x) = -x\) and \(\mathrm{Tr}_{\mathrm{D}/k_1}(x) = 0\) . Conversely, if \(\Delta\) is a central simple \(k_{1}\) -algebra of dimension \((l + 1)^2\) , equipped with an involutive anti-automorphism \(\sigma\) whose restriction to \(k_{1}\) is \(x \mapsto \bar{x}\) , the Lie algebra \(\mathfrak{s}\mathfrak{u}_{\Delta,\sigma}\) is absolutely simple of type \(\mathrm{A}_l\) , and the corresponding homomorphism \(\pi\) is that associated to \(k_{1}\) . Two such algebras \(\mathfrak{s}\mathfrak{u}_{\Delta,\sigma}\) and \(\mathfrak{s}\mathfrak{u}_{\Delta',\sigma'}\) are isomorphic if and only if there exists a \(k\) -isomorphism \(f: \Delta \to \Delta'\) such that \(\sigma' \circ f = f \circ \sigma\) .

Show that, when \(\mathrm{D} = \mathbf{M}_{l+1}(k_{1})\) , there exists an invertible hermitian matrix H of degree \(l + 1\) , unique up to multiplication by an element of \(k^{*}\) , such that

\[
\sigma (x) = H. ^ {t} \bar {x}. H ^ {- 1}
\]

for all \(x \in \mathbf{M}_{l+1}(k_1)\) ; the algebra g can then be identified with the algebra \(\mathfrak{su}(l + 1, H)\) consisting of the matrices x such that \(x.H + H.^t\bar{x} = 0\) and \(\operatorname{Tr}(x) = 0\) .

¶ 17) Let g be an absolutely simple Lie algebra of type \(B_{l}\) (resp. \(C_{l}, D_{l}\) ), with \(l \geq 2\) (resp. \(l \geq 3, l \geq 4\) ). Assume further that, when g is of type \(D_{4}\) , the image of the homomorphism \(\pi : \operatorname{Gal}(\bar{k}/k) \to \operatorname{Aut}(\bar{\mathrm{R}}, \bar{\mathrm{B}})\) defined in Exerc. 8 of §5 is of order \(\leq 2\) . Show that there exists a central simple algebra D of dimension \((2l + 1)^{2}\) (resp. \(4l^{2}, 4l^{2}\) ) and an involutive anti-automorphism \(\sigma\) of D such that g is isomorphic to the Lie subalgebra of D consisting of the elements x such that \(\sigma(x) = -x\) and \(\operatorname{Tr}_{\mathrm{D}}(x) = 0\) (same method as in Exerc. 16 a)). \(^{22}\)

\begin{enumerate}
\def\labelenumi{\arabic{enumi})}
\setcounter{enumi}{17}
\tightlist
\item
  Let V be a finite dimensional vector space, Q a non-degenerate quadratic form on V, and \(\Psi\) the symmetric bilinear form associated to Q. Denote the Lie algebra \(\mathfrak{o}(\varPsi)\) by g and the extension of \(\varPsi\) to \(\bigwedge^{2}(V)\) by \(\varPsi_{2}\) .
\end{enumerate}

\begin{enumerate}
\def\labelenumi{\alph{enumi})}
\tightlist
\item
  Show that there exists an isomorphism of vector spaces \(\theta : \bigwedge^{2}(\mathrm{V}) \to \mathfrak{g}\) characterized by the following equivalent properties:
\end{enumerate}

\begin{enumerate}
\def\labelenumi{(\roman{enumi})}
\item
  For \(a, b\) and \(x\) in \(\mathrm{V}\) , we have \(\theta(a \wedge b). x = a. \Psi(x, b) - b. \Psi(x, a)\) .
\item
  For \(x, y\) in V and \(u\) in \(\bigwedge^2 (\mathrm{V})\) we have \(\varPsi_2(x\wedge y,u) = \varPsi (x,\theta (u).y)\) .
\end{enumerate}

Let \(\sigma\) be the identity representation of \(\mathfrak{g}\) on \(V\) ; then \(\theta\) is an isomorphism of \(\bigwedge^{2}(V)\) with the adjoint representation of \(\mathfrak{g}\) .

\begin{enumerate}
\def\labelenumi{\alph{enumi})}
\setcounter{enumi}{1}
\item
  Define the linear map \(f: \mathfrak{g} \to \mathrm{C}^{+}(\mathrm{Q})\) as in Lemma 1 of no. 2. Show that \(f\theta(a \wedge b) = \frac{1}{2}(ab - ba)\) for a, b in V, and deduce a new proof of assertions (iii), (iv) and (v) of Lemma 1 of no. 2.
\item
  The notations \(l, \mathrm{F}, \mathrm{F}', e_0, \mathrm{N}, \lambda\) and \(\rho\) are those of no. 2. Choose \(e \neq 0\) in \(\bigwedge (\mathrm{F}')\) and define the bilinear form \(\varPhi\) on \(\mathrm{N}\) by
\end{enumerate}

\[
x \wedge y = (- 1) ^ {p (p + 1) / 2} \Phi (x, y). e \quad (x \in \bigwedge^ {p} (\mathrm{F} ^ {\prime}), y \in \mathrm{N}).
\]

Show that \(\Phi(\lambda(a).x,y)+\Phi(x,\lambda(a).y)=0\) for x,y in N and a in \(F\oplus F'\) . Deduce that \(\Phi\) is invariant under the representation \(\rho\) of g (remark that the Lie algebra \(\rho(\mathfrak{g})\) is generated by \(\lambda(F\oplus F')\) by a) and b) above).

\begin{enumerate}
\def\labelenumi{\arabic{enumi})}
\setcounter{enumi}{18}
\tightlist
\item
  Let \(\mathrm{V}\) be a finite dimensional vector space and \(\mathrm{E} = \mathrm{V} \oplus \mathrm{V}^*\) . Define a non-degenerate bilinear form \(\varPhi\) on \(\mathrm{E}\) by
\end{enumerate}

\[
\Phi ((x, x ^ {*}), (y, y ^ {*})) = \langle x, y ^ {*} \rangle + \langle y, x ^ {*} \rangle .
\]

Put \(N = \bigwedge(V)\) and \(Q(x) = \frac{1}{2}\Phi(x, x)\) for \(x \in E\) . Denote the spinor representation by \(\lambda : C(Q) \to \text{End}(N)\) as in no. 2.IV, define \(f : \mathfrak{o}(\Phi) \to C^{+}(Q)\) as in Lemma 1 of no. 2, and denote by \(\rho\) the linear representation \(\lambda \circ f\) of the algebra \(\mathfrak{o}(\Phi)\) on N.

\begin{enumerate}
\def\labelenumi{\alph{enumi})}
\item
  Associate to any endomorphism \(u\) of \(V\) the endomorphism \(\tilde{u}\) of \(E\) by the formula \(\tilde{u}(x, x^{*}) = (u(x), -^{t}u(x^{*}))\) . Show that \(u \mapsto \tilde{u}\) is a homomorphism of Lie algebras from \(\mathfrak{gl}(V)\) to \(\mathfrak{o}(\varPhi)\) ; moreover, for \(u\) in \(\mathfrak{gl}(V)\) , \(\rho(\tilde{u})\) is the unique derivation of the algebra \(\bigwedge(V) = N\) that coincides with \(u\) on \(V\) .
\item
  Let \(\varPsi\) be a non-degenerate alternating bilinear form on V and \(\gamma: V \to V^{*}\) the isomorphism defined by \(\varPsi(x, y) = \langle x, \gamma(y) \rangle\) for x, y in V. Show that the endomorphisms \(\tilde{X}_{+}\) and \(\tilde{X}_{-}\) of E defined by
\end{enumerate}

\[
\tilde {X} _ {+} (x, x ^ {*}) = (\gamma^ {- 1} (x ^ {*}), 0), \quad \tilde {X} _ {-} (x, x ^ {*}) = (0, - \gamma (x))
\]

belongs to \(\mathfrak{o}(\varPhi)\) . Put \(\tilde{H} = (-1)^{\sim}\) . Show that \((\tilde{H}, \tilde{X}_{+}, \tilde{X}_{-})\) is an \(\mathfrak{sl}_2\) -triplet in the Lie algebra \(\mathfrak{o}(\varPhi)\) .

\begin{enumerate}
\def\labelenumi{\alph{enumi})}
\setcounter{enumi}{2}
\tightlist
\item
  Show that \(\rho\) takes the endomorphisms \(\tilde{H}, \tilde{X}_{+}, \tilde{X}_{-}\) of \(\mathfrak{o}(\Phi)\) to the endomorphisms of N denoted by \(H, X_{+}\) and \(X_{-}\) , respectively, in no. 3.IV. Deduce that \((H, X_{+}, X_{-})\) is an \(\mathfrak{sl}_{2}\) -triplet in the Lie algebra \(\mathfrak{gl}(N)\) .
\end{enumerate}

